% SPDX-License-Identifier: CC-BY-SA-4.0
% Copyright (c) 2026 Rafael Beraldo
% Exemplos de cada tipo de slide, herdados do modelo. Não entram na
% apresentação (veja presentation.tex); apague quando não forem mais úteis.
% Incluído por presentation.tex; compile a partir dele.

% --- Lista com revelação progressiva ----------------------------------------
\begin{frame}{Quando o celular não ajuda}
  \begin{itemize}[<+->]
    \item Sem cobertura em vales, encostas e trilhas longas
    \item Bateria acaba justamente quando mais precisamos
    \item O rádio fala com \alert{todo mundo ao alcance}, não com um número
    \item Um rádio de \alert{VHF} de mão custa menos que uma bota
  \end{itemize}
\end{frame}

% --- Duas colunas: texto + imagem -------------------------------------------
\begin{frame}{Rádio e montanhismo se complementam}
  \begin{columns}[T,onlytextwidth]
    \column{0.55\textwidth}
      \textbf{O que o rádio resolve}
      \begin{itemize}
        \item Coordenação entre cordadas
        \item Contato com a base e com equipes de resgate
        \item Boletim meteorológico em tempo real
      \end{itemize}
      \vspace{0.5em}
      \textbf{O que o rádio exige}
      \begin{itemize}
        \item Licença da Anatel e indicativo
        \item Disciplina de operação
      \end{itemize}
    \column{0.4\textwidth}
      % Troque example-image-a pelo nome da sua foto em img/
      \foto{example-image-a}{Legenda da foto}
  \end{columns}
\end{frame}

% --- Foto em tela cheia com faixa de legenda --------------------------------
% \cerjfullimage[opções]{arquivo}{legenda}; arquivo em img/
\cerjfullimage{example-image-16x9}{%
  \usebeamerfont{frametitle}Pedra do Sino, Teresópolis\par
  \small Repetidora alcançável a partir do cume}

% =============================================================================
\section{Conceitos básicos}
% =============================================================================

% --- Tabela -----------------------------------------------------------------
\begin{frame}{Faixas mais usadas pelo montanhista}
  % Valores ilustrativos: confira no plano de bandas antes de apresentar
  \begin{center}
    \begin{tabular}{@{}llll@{}}
      \toprule
      Banda & Faixa & Alcance típico & Uso \\
      \midrule
      2 m   & \freqrange{144}{148}   & Linha de visada & Simplex local, repetidoras \\
      70 cm & \freqrange{430}{440}   & Linha de visada & Simplex, repetidoras urbanas \\
      40 m  & \freqrange{7,0}{7,3}   & Centenas de km   & Contato com a base distante \\
      \bottomrule
    \end{tabular}
  \end{center}
  \vspace{0.5em}
  Frequência de chamada simplex em VHF: \freq{146,520}
\end{frame}

% --- Blocos do beamer -------------------------------------------------------
\begin{frame}{Blocos: definição, alerta e exemplo}
  \begin{block}{Simplex}
    Comunicação direta entre dois rádios, na mesma frequência, sem repetidora.
  \end{block}
  \begin{alertblock}{Emergência}
    Em emergência real, qualquer frequência pode ser usada para pedir socorro.
  \end{alertblock}
  \begin{exampleblock}{Exemplo de chamada}
    “\indicativo{PY1XXX} chamando \indicativo{PY1YYY}, câmbio.”
  \end{exampleblock}
\end{frame}

% --- Caixas de destaque (tcolorbox) -----------------------------------------
\begin{frame}{Caixas de destaque}
  \begin{dica}
    Leve uma antena telescópica de meia onda: dobra o alcance do rádio de mão.
  \end{dica}
  \begin{atencao}
    Não transmita com a antena desconectada, o estágio final do rádio queima.
  \end{atencao}
  \begin{regulamento}
    Operar sem licença é infração à Lei Geral de Telecomunicações.
  \end{regulamento}
\end{frame}

% --- Lista numerada: passo a passo ------------------------------------------
\begin{frame}{Como pedir ajuda pelo rádio}
  \begin{enumerate}
    \item Respire, escute a frequência por alguns segundos
    \item Diga seu indicativo e “\alert{emergência}” três vezes
    \item Informe \alert{onde}, \alert{quem} e \alert{o quê}
    \item Fique na escuta, não desligue o rádio
  \end{enumerate}
\end{frame}

% --- Frase de impacto (standout do metropolis) ------------------------------
\begin{frame}[standout]
  Sem sinal de celular? \\ Sem problema.
\end{frame}

%%% Local Variables:
%%% TeX-master: "../presentation"
%%% fill-column: 100
%%% End:
